\scanpage{28}

$P^1_S$ sections $0,1,\infty$. Le groupe des automorphismes s'identifie au groupe $\mathfrak S_3$ des permutations des trois \unsure{pts} $0,1,\infty$, groupe à $6$ éléments $1,\varphi,\psi,\lambda,\mu,\nu$.
\[
(1)\quad\left\{\begin{array}{l}\varphi(t)=1-\dfrac1t\\\psi(t)=\dfrac1{1-t}\end{array}\right.
\qquad\text{2 permutations circulaires non triviales de }0,1,\infty
\]
\[
(2)\quad\left\{\begin{array}{ll||l}
\lambda t=\dfrac1t&\text{transposition de }0,\infty&1\\
\mu t=1-t&\text{transposition de }0,1&\infty\\
\nu t=1+\dfrac1{t-1}=\dfrac t{t-1}&\text{transposition de }1,\infty&0
\end{array}\right.
\]
\margin{\unsure{renommer}\quad
\[\left\{\begin{array}{l}\lambda t=1-\unsure{t}\\\mu t=\dfrac t{\unsure{t-1}}\\\nu t=\unsure{1/t}\\?\ ?\end{array}\right.\]}
\[
3\ a)\quad\boxed{\begin{array}{ll}\varphi^2=\psi&\psi^2=\varphi\\\deleted{\unsure{\psi^2=\varphi|\varphi^2=\psi}}&\varphi\psi=\psi\varphi=1\end{array}}
\]
\[
b)\quad\boxed{\lambda^2=\mu^2=\nu^2=1}
\]
\[
c)\quad\boxed{\begin{array}{l}\lambda\mu=\mu\nu=\nu\lambda=\varphi\\\mu\lambda=\nu\mu=\lambda\nu=\psi\end{array}}
\]
\diagram{\begin{array}{c}\displaystyle d)\quad\boxed{\begin{array}{l|l}\varphi\lambda=\mu\varphi=\nu&\psi\nu=\mu\psi=\lambda\\\varphi\mu=\nu\varphi=\lambda&\psi\mu=\lambda\psi=\nu\\\varphi\nu=\lambda\varphi=\mu&\psi\lambda=\nu\psi=\mu\end{array}}\end{array}}

$1,\varphi,\psi$ \unsure{gr. gr. cyclique} d'ordre $3$, de générateurs $\varphi$ et $\psi$, invariant dans $\mathfrak S_3$. Le quotient est $\mathbb Z/2$ et $\mathfrak S_3$ est produit semi-direct ; il y a trois sous-groupes sections, engendrés resp\unsure{\textsuperscript{t}} par $\lambda,\mu,\nu$ (leurs seuls éléments d'ordre $2$)
\[
4)\quad\left\{\begin{array}{lll}\varphi_\lambda=\mu&\varphi_\mu=\nu&\varphi_\nu=\lambda\\\psi_\nu=\mu&\psi_\mu=\lambda&\psi_\lambda=\nu\end{array}\right.
\]
\[
5)\quad\left\{\begin{array}{l}\lambda_\varphi=\mu_\varphi=\nu_\varphi=\psi\\\lambda_\psi=\mu_\psi=\nu_\psi=\varphi\end{array}\right.
\]
De plus, on a dans $P^1$ un automorphisme $\sigma$ \deleted{\unsure{de}} de carré $1$, donné par
\[
6)\quad\sigma t=\frac{t-2}{2t-1}\qquad[\sigma^2=\mathrm{id}]
\]
\unsure{permutant à} : $\unsure{\mathfrak S_6|\mathfrak S_3}$ (c'est d'ailleurs le seul, \unsure{sauf|avec}
\mspage{27}
